% !TeX root = ../main.tex
%%% 多小问大题：逐问作答流示范
\begin{problem}[20]{实对称矩阵的特征值与正定性}
\label{sec:matrix}
设实二次型 $f(x_1, x_2, x_3) = x_1^2 + x_2^2 + x_3^2 + 2x_1 x_2 + 2x_2 x_3 + 2x_3 x_1$ 对应的实对称矩阵为 $A$。
\end{problem}

\begin{subproblems}
  \subproblem[8]{%
    \label{sub:eigen}%
    求矩阵 $A$ 的全部特征值，并写出对应特征子空间的一组正交基。
  }
  \begin{solution}
    易写出二次型对应的实对称矩阵：
    \[
      A = \begin{pmatrix} 1 & 1 & 1 \\ 1 & 1 & 1 \\ 1 & 1 & 1 \end{pmatrix}.
    \]

    计算特征多项式得 $\det(\lambda I - A) = \lambda^2(\lambda - 3)$，因此特征值为 $\lambda_1 = 3, \lambda_{2,3} = 0$。
    
    对于 $\lambda = 3$，特征子空间基向量为 $(1, 1, 1)^\top$，单位化得：
    \[
      \vect{q}_1 = \frac{1}{\sqrt{3}}(1, 1, 1)^\top.
    \]

    对于重特征值 $\lambda = 0$，解方程 $x_1 + x_2 + x_3 = 0$，施密特正交化并单位化得：
    \[
      \vect{q}_2 = \frac{1}{\sqrt{2}}(1, -1, 0)^\top, \quad
      \vect{q}_3 = \frac{1}{\sqrt{6}}(1, 1, -2)^\top.
    \]
  \end{solution}

  \subproblem[12]{%
    \label{sub:posdef}%
    判断矩阵 $A$ 是否正定，并求正交矩阵 $P$ 使得 $P^\top A P$ 为对角阵。
  }
  \begin{solution}
    矩阵 $A$ 不是正定矩阵，而是半正定的。
    因为若实对称阵正定，则其全部特征值必严格大于 $0$；而由 \cref{sub:eigen} 可知 $\lambda_2 = \lambda_3 = 0 \ngtr 0$，故 $A$ 不是正定矩阵。

    取正交矩阵 $P = (\vect{q}_1, \vect{q}_2, \vect{q}_3)$，由 \cref{sub:eigen} 求得的正交基可得：
    \[
      P^\top A P = \operatorname{diag}(3, 0, 0).
    \]
  \end{solution}
\end{subproblems}
